\documentclass[11pt,a4paper]{article}
\usepackage[utf8]{inputenc}
\usepackage[italian]{babel}
\usepackage{geometry}
\usepackage{graphicx}
\usepackage{hyperref}
\usepackage{listings}
\usepackage{xcolor}
\usepackage{amsmath}
\usepackage{enumitem}
\usepackage{booktabs}
\usepackage{longtable}
\usepackage{array}
\usepackage{tcolorbox}

\geometry{margin=2.5cm}

% Configurazione codice
\lstset{
    language=JavaScript,
    basicstyle=\ttfamily\small,
    keywordstyle=\color{blue}\bfseries,
    stringstyle=\color{red},
    commentstyle=\color{green!60!black},
    numberstyle=\tiny\color{gray},
    numbers=left,
    numbersep=5pt,
    frame=single,
    breaklines=true,
    showstringspaces=false,
    tabsize=2,
    escapeinside={(*@}{@*)}
}

% Box per evidenziare sezioni importanti
\newtcolorbox{importantbox}[1]{
    colback=red!5!white,
    colframe=red!75!black,
    title=#1,
    fonttitle=\bfseries
}

\newtcolorbox{warningbox}[1]{
    colback=orange!5!white,
    colframe=orange!75!black,
    title=#1,
    fonttitle=\bfseries
}

\newtcolorbox{tipbox}[1]{
    colback=blue!5!white,
    colframe=blue!75!black,
    title=#1,
    fonttitle=\bfseries
}

\title{\textbf{Analisi Componente: TenantContext}\\
\large Documentazione Tecnica Esaustiva}
\author{Documentazione Sistema}
\date{\today}

\begin{document}

\maketitle
\tableofcontents
\newpage

\section{Storico Modifiche}

\begin{importantbox}{Correzioni Applicate - \today}
Le seguenti vulnerabilità e miglioramenti sono stati implementati:
\begin{itemize}
    \item \textbf{Error Handling Standardizzato}: Uso di \texttt{AppError.unauthorized()} e \texttt{AppError.validation()} invece di \texttt{res.status().json()} manuale
    \item \textbf{Validazione ObjectId Rigorosa}: Controllo con \texttt{mongoose.Types.ObjectId.isValid()} prima della conversione - \textbf{Previene crash con ID malformati}
    \item \textbf{Validazione Stringa Vuota}: Check esplicito per \texttt{req.user.id === ''} - \textbf{Previene ObjectId("") che crea ID invalido}
    \item \textbf{Object.freeze() su req.tenantScope}: Prevenzione modifiche accidentali da middleware successivi - \textbf{Previene bug di sviluppo}
\end{itemize}
\end{importantbox}

\section{Responsabilità}

\subsection{Cosa fa esattamente questo componente?}

\textbf{TenantContext} è un middleware Express che crea un "ponte" tra il layer di autenticazione (JWT) e il layer di accesso ai dati (BaseService/Mongoose), iniettando il contesto tenant nella request.

\textbf{Funzioni principali}:
\begin{itemize}
    \item \textbf{Estrazione Identità Utente}: Legge \texttt{req.user.id} dal middleware di autenticazione
    \item \textbf{Conversione ObjectId}: Converte string ID in \texttt{mongoose.Types.ObjectId} se necessario
    \item \textbf{Iniezione Context}: Crea \texttt{req.tenantScope} con metodi helper per operazioni tenant-scoped
    \item \textbf{Gestione Edge Cases}: Gestisce richieste non autenticate (opzionale o errore 401)
    \item \textbf{Validazione Plugin}: Verifica che i modelli abbiano il plugin multi-tenancy applicato
\end{itemize}

\subsection{Perché esiste?}

\begin{importantbox}{Problema Risolto}
Senza TenantContext, ogni service dovrebbe:
\begin{itemize}
    \item Estrarre manualmente \texttt{req.user.id} in ogni controller
    \item Convertire string ID in ObjectId ogni volta
    \item Gestire edge cases (utente non autenticato, ID malformato)
    \item Verificare che i modelli abbiano il plugin applicato
    \item Scrivere codice duplicato per creare filtri tenant
\end{itemize}
\end{importantbox}

\textbf{Vantaggi dell'approccio}:
\begin{enumerate}
    \item \textbf{Separation of Concerns}: Autenticazione separata da logica business
    \item \textbf{DRY (Don't Repeat Yourself)}: Logica di estrazione centralizzata
    \item \textbf{Type Safety}: Conversione ObjectId centralizzata e consistente
    \item \textbf{API Helper}: Metodi \texttt{model()}, \texttt{create()}, \texttt{owns()} rendono il codice più leggibile
    \item \textbf{Fail-Fast}: Errore esplicito se utente non autenticato (quando required)
\end{enumerate}

\section{Struttura del Codice}

\subsection{Architettura}

TenantContext implementa il pattern \textbf{Middleware Chain} di Express:

\begin{lstlisting}
[HTTP Request]
     │
     ▼
[Auth Middleware] ──► req.user = { id, email }
     │
     ▼
[Tenant Middleware] ──► req.tenantScope = { tenantId, model(), create(), owns(), filter }
     │
     ▼
[Route Handler] ──► Usa req.tenantScope
     │
     ▼
[Service Layer] ──► BaseService usa req.tenantScope.tenantId
     │
     ▼
[Mongoose Query] ──► Filtro automatico { user: tenantId }
\end{lstlisting}

\subsection{Metodi Principali}

\subsubsection{tenantContext(options)}

\begin{lstlisting}
const tenantContext = (options = {}) => {
    const { required = true } = options;

    return (req, res, next) => {
        // 1. Validazione: Utente presente?
        if (!req.user || !req.user.id || req.user.id === '') {
            if (required) {
                // FIX: Uso di AppError standardizzato
                return next(AppError.unauthorized(
                    'Tenant context required. Please authenticate.',
                    {
                        code: 'TENANT_CONTEXT_MISSING',
                        category: 'AUTHORIZATION',
                        suggestion: 'Assicurati di essere autenticato...',
                    }
                ));
            }
            return next();
        }

        // 2. Validazione: ID valido?
        // FIX: Controllo rigoroso dell'ID
        let tenantId;
        const userId = req.user.id;

        if (userId instanceof mongoose.Types.ObjectId) {
            tenantId = userId;
        } else if (typeof userId === 'string' && mongoose.Types.ObjectId.isValid(userId)) {
            tenantId = new mongoose.Types.ObjectId(userId);
        } else {
            return next(AppError.validation(
                'Invalid user ID format in token.',
                {},
                { code: 'INVALID_USER_ID', category: 'VALIDATION' }
            ));
        }

        // 3. Crea req.tenantScope con helper methods
        // FIX: Object.freeze() per prevenire modifiche accidentali
        req.tenantScope = Object.freeze({
            tenantId,
            model: (Model) => { /* ... */ },
            create: async (Model, data) => { /* ... */ },
            owns: async (Model, documentId) => { /* ... */ },
            get filter() { return { user: tenantId }; },
        });

        next();
    };
};
\end{lstlisting}

\textbf{Logica}:
\begin{itemize}
    \item \textbf{Estrazione Utente}: Legge \texttt{req.user.id} (popolato da \texttt{requireAuth} middleware)
    \item \textbf{Edge Case Handling}: Se \texttt{required=true} e utente mancante, ritorna 401 immediatamente
    \item \textbf{Conversione Type}: Converte string ID in ObjectId per compatibilità MongoDB
    \item \textbf{Iniezione Context}: Crea \texttt{req.tenantScope} con metodi helper
\end{itemize}

\subsubsection{req.tenantScope.model(Model)}

\begin{lstlisting}
model: (Model) => {
    if (!Model.forTenant) {
        throw new Error(
            `Model ${Model.modelName} does not have multi-tenancy plugin. ` +
            'Did you forget to apply the plugin?'
        );
    }
    return Model.forTenant(tenantId);
}
\end{lstlisting}

\textbf{Logica}:
\begin{itemize}
    \item \textbf{Validazione Plugin}: Verifica che il Model abbia il metodo \texttt{forTenant} (aggiunto dal plugin)
    \item \textbf{Scoping Automatico}: Ritorna un Proxy del Model con contesto tenant iniettato
    \item \textbf{Esempio}: \texttt{req.tenantScope.model(Project).find()} esegue automaticamente \texttt{Project.find(\{user: tenantId\})}
\end{itemize}

\subsubsection{req.tenantScope.create(Model, data)}

\begin{lstlisting}
create: async (Model, data) => {
    if (!Model.createForTenant) {
        throw new Error(
            `Model ${Model.modelName} does not have multi-tenancy plugin.`
        );
    }
    return Model.createForTenant(tenantId, data);
}
\end{lstlisting}

\textbf{Logica}:
\begin{itemize}
    \item \textbf{Validazione Plugin}: Verifica che il Model abbia il metodo \texttt{createForTenant}
    \item \textbf{Creazione Sicura}: Inietta \texttt{tenantId} nel documento prima della creazione
    \item \textbf{Prevenzione Mass Assignment}: Il campo \texttt{user} viene iniettato server-side, non dal client
\end{itemize}

\subsubsection{req.tenantScope.owns(Model, documentId)}

\begin{lstlisting}
owns: async (Model, documentId) => {
    if (!Model.belongsToTenant) {
        throw new Error(
            `Model ${Model.modelName} does not have multi-tenancy plugin.`
        );
    }
    return Model.belongsToTenant(tenantId, documentId);
}
\end{lstlisting}

\textbf{Logica}:
\begin{itemize}
    \item \textbf{Verifica Ownership}: Controlla se un documento appartiene al tenant corrente
    \item \textbf{Validazione Esplicita}: Utile per validazioni esplicite prima di operazioni sensibili
    \item \textbf{Esempio}: \texttt{if (!await req.tenantScope.owns(Project, id)) throw AppError.forbidden()}
\end{itemize}

\subsubsection{req.tenantScope.filter}

\begin{lstlisting}
get filter() {
    return { user: tenantId };
}
\end{lstlisting}

\textbf{Logica}:
\begin{itemize}
    \item \textbf{Filtro Base}: Ritorna il filtro base per query manuali
    \item \textbf{Query Complesse}: Utile per aggregation pipeline o query non coperte dagli helper
    \item \textbf{Esempio}: \texttt{Project.aggregate([\{ \$match: \{ ...req.tenantScope.filter, status: 'active' \} \}])}
\end{itemize}

\subsubsection{withTenantAuth(authMiddleware)}

\begin{lstlisting}
const withTenantAuth = (authMiddleware) => {
    return [authMiddleware, tenantContext({ required: true })];
};
\end{lstlisting}

\textbf{Logica}:
\begin{itemize}
    \item \textbf{Convenienza}: Combina auth + tenant context in un'unica chiamata
    \item \textbf{Utilizzo}: \texttt{router.use(...withTenantAuth(requireAuth))}
    \item \textbf{DRY}: Evita di scrivere \texttt{[requireAuth, tenantContext(\{required: true\})]} ogni volta
\end{itemize}

\section{Dipendenze}

\subsection{Dipendenze Dirette}

\begin{itemize}
    \item \textbf{mongoose}: Per \texttt{mongoose.Types.ObjectId} (conversione ID)
    \item \textbf{req.user}: Popolato dal middleware \texttt{requireAuth} (deve essere applicato PRIMA)
\end{itemize}

\subsection{Dipendenze Indirette}

\begin{itemize}
    \item \textbf{requireAuth Middleware}: Estrae token JWT dal cookie e popola \texttt{req.user}
    \item \textbf{multiTenancyPlugin}: I modelli devono avere il plugin applicato per usare \texttt{model()}, \texttt{create()}, \texttt{owns()}
\end{itemize}

\subsection{Ordine Middleware}

\begin{warningbox}{⚠️ Ordine Critico}
Il middleware \texttt{tenantContext} DEVE essere applicato DOPO \texttt{requireAuth}:

\begin{lstlisting}
// ✅ CORRETTO
router.use(requireAuth);
router.use(tenantContext({ required: true }));

// ❌ SBAGLIATO - req.user non esiste ancora
router.use(tenantContext({ required: true }));
router.use(requireAuth);
\end{lstlisting}
\end{warningbox}

\section{Analisi Critica}

\subsection{Punti di Forza}

\begin{enumerate}
    \item \textbf{Fail-Fast su Utente Mancante}:
    \begin{itemize}
        \item Se \texttt{required=true} e \texttt{req.user.id} manca, ritorna 401 immediatamente
        \item Previene esecuzione di codice con contesto invalido
        \item Messaggio di errore chiaro: \texttt{TENANT\_CONTEXT\_MISSING}
    \end{itemize}
    
    \item \textbf{Conversione Type Sicura}:
    \begin{itemize}
        \item Gestisce sia string ID che ObjectId
        \item Previene errori MongoDB per ID malformati
        \item Conversione centralizzata (DRY)
    \end{itemize}
    
    \item \textbf{API Helper Elegante}:
    \begin{itemize}
        \item \texttt{model()}, \texttt{create()}, \texttt{owns()} rendono il codice più leggibile
        \item Validazione plugin automatica (fail-fast se plugin mancante)
        \item Filtro base accessibile tramite getter
    \end{itemize}
    
    \item \textbf{Modo Opzionale}:
    \begin{itemize}
        \item \texttt{tenantContext(\{required: false\})} permette routes pubbliche
        \item Se utente non autenticato, continua senza \texttt{req.tenantScope}
        \item Utile per routes che funzionano sia per utenti autenticati che non
    \end{itemize}
    
    \item \textbf{withTenantAuth Helper}:
    \begin{itemize}
        \item Riduce boilerplate code
        \item Garantisce ordine corretto dei middleware
    \end{itemize}
\end{enumerate}

\subsection{Possibili Bug/Rischi}

\begin{enumerate}
    \item \textbf{req.user.id è undefined/null ma non string/ObjectId}:
    \begin{itemize}
        \item \textbf{Stato attuale}: ✅ \textbf{RISOLTO} - Check esplicito \texttt{!req.user || !req.user.id || req.user.id === ''}
        \item \textbf{Implementazione}: Validazione fail-fast con \texttt{AppError.unauthorized()} standardizzato
        \item \textbf{Protezione}: Previene crash con stringa vuota che passerebbe il check \texttt{!req.user?.id}
    \end{itemize}
    
    \item \textbf{ObjectId malformato}:
    \begin{itemize}
        \item \textbf{Stato attuale}: ✅ \textbf{RISOLTO} - Validazione con \texttt{mongoose.Types.ObjectId.isValid()} prima della conversione
        \item \textbf{Implementazione}: 
        \begin{itemize}
            \item Se \texttt{userId instanceof mongoose.Types.ObjectId} → usa direttamente
            \item Se \texttt{typeof userId === 'string' && mongoose.Types.ObjectId.isValid(userId)} → converte
            \item Altrimenti → errore 400 con \texttt{AppError.validation()}
        \end{itemize}
        \item \textbf{Protezione}: Previene crash MongoDB e creazione di ObjectId invalidi
    \end{itemize}
    
    \item \textbf{req.user.id è ObjectId ma di un altro utente}:
    \begin{itemize}
        \item \textbf{Stato attuale}: Non c'è validazione che \texttt{req.user.id} corrisponda all'utente autenticato
        \item \textbf{Rischio}: Se il token JWT è compromesso o modificato, potrebbe iniettare un ID diverso
        \item \textbf{Mitigazione}: Il middleware \texttt{requireAuth} verifica la firma JWT, quindi questo è già protetto
        \item \textbf{Miglioramento Futuro}: Considerare verifica aggiuntiva che l'utente esista nel DB (ma aumenta latenza)
    \end{itemize}
    
    \item \textbf{Plugin mancante su Model}:
    \begin{itemize}
        \item \textbf{Stato attuale}: ✅ \textbf{GESTITO} - \texttt{model()}, \texttt{create()}, \texttt{owns()} lanciano errore esplicito se plugin mancante
        \item \textbf{Protezione}: Fail-fast con messaggio chiaro: \texttt{"Model X does not have multi-tenancy plugin"}
    \end{itemize}
    
    \item \textbf{req.tenantScope sovrascritto}:
    \begin{itemize}
        \item \textbf{Stato attuale}: ✅ \textbf{RISOLTO} - \texttt{Object.freeze()} applicato su \texttt{req.tenantScope}
        \item \textbf{Implementazione}: \texttt{req.tenantScope = Object.freeze(\{ ... \})}
        \item \textbf{Protezione}: Previene modifiche accidentali da middleware successivi (fail-fast con TypeError)
    \end{itemize}
    
    \item \textbf{Concorrenza su req.tenantScope}:
    \begin{itemize}
        \item \textbf{Stato attuale}: Ogni request ha il proprio \texttt{req.tenantScope}
        \item \textbf{Rischio}: Nessuno (Express è single-threaded per request)
    \end{itemize}
    
    \item \textbf{Memory Leak}:
    \begin{itemize}
        \item \textbf{Stato attuale}: \texttt{req.tenantScope} viene creato per ogni request
        \item \textbf{Rischio}: Nessuno (viene garbage collected dopo la response)
    \end{itemize}
\end{enumerate}

\subsection{Performance}

\begin{enumerate}
    \item \textbf{Conversione ObjectId}:
    \begin{itemize}
        \item \textbf{Operazione}: \texttt{new mongoose.Types.ObjectId(stringId)}
        \item \textbf{Costo}: O(1), operazione molto veloce
        \item \textbf{Impatto}: Trascurabile (microsecondi)
    \end{itemize}
    
    \item \textbf{Creazione req.tenantScope}:
    \begin{itemize}
        \item \textbf{Operazione}: Creazione oggetto con 4 metodi
        \item \textbf{Costo}: O(1), allocazione memoria minima
        \item \textbf{Impatto}: Trascurabile
    \end{itemize}
    
    \item \textbf{Validazione Plugin}:
    \begin{itemize}
        \item \textbf{Operazione}: Check \texttt{Model.forTenant} (property lookup)
        \item \textbf{Costo}: O(1)
        \item \textbf{Impatto}: Trascurabile, ma viene eseguito ad ogni chiamata \texttt{model()}
        \item \textbf{Miglioramento Futuro}: Cache del risultato per Model (ma probabilmente over-engineering)
    \end{itemize}
\end{enumerate}

\subsection{Miglioramenti Codice (Clean Code)}

\begin{tipbox}{Suggerimenti per Migliorare}
\end{tipbox}

\begin{enumerate}
    \item \textbf{Validazione ObjectId}:
    \begin{lstlisting}
// ✅ IMPLEMENTATO - Validazione rigorosa con AppError
let tenantId;
const userId = req.user.id;

if (userId instanceof mongoose.Types.ObjectId) {
    tenantId = userId;
} else if (typeof userId === 'string' && mongoose.Types.ObjectId.isValid(userId)) {
    tenantId = new mongoose.Types.ObjectId(userId);
} else {
    // Se l'ID è spazzatura (es. stringa vuota o malformata), errore 400
    return next(AppError.validation(
        'Invalid user ID format in token.',
        {},
        {
            code: 'INVALID_USER_ID',
            category: 'VALIDATION',
            suggestion: 'Il token di autenticazione contiene un ID utente non valido. Effettua nuovamente il login.',
        }
    ));
}
    \end{lstlisting}
    
    \item \textbf{Validazione Stringa Vuota}:
    \begin{lstlisting}
// ✅ IMPLEMENTATO - Check esplicito per stringa vuota
if (!req.user || !req.user.id || req.user.id === '') {
    if (required) {
        return next(AppError.unauthorized(
            'Tenant context required. Please authenticate.',
            {
                code: 'TENANT_CONTEXT_MISSING',
                category: 'AUTHORIZATION',
                suggestion: 'Assicurati di essere autenticato e che il middleware requireAuth sia applicato prima di tenantContext.',
            }
        ));
    }
    return next();
}
    \end{lstlisting}
    
    \item \textbf{Freeze req.tenantScope}:
    \begin{lstlisting}
// ✅ IMPLEMENTATO - Prevenire modifiche accidentali
req.tenantScope = Object.freeze({
    tenantId,
    model: (Model) => { /* ... */ },
    create: async (Model, data) => { /* ... */ },
    owns: async (Model, documentId) => { /* ... */ },
    get filter() { return { user: tenantId }; },
});
    \end{lstlisting}
    
    \item \textbf{TypeScript per Type Safety}:
    \begin{itemize}
        \item Convertire in TypeScript per type checking
        \item Definire interfaccia \texttt{ITenantScope} per autocomplete
        \item Previene errori a compile-time
    \end{itemize}
    
    \item \textbf{Logging per Debug}:
    \begin{lstlisting}
// Aggiungere logging opzionale
if (process.env.DEBUG_TENANT === 'true') {
    console.log(`[TenantContext] User ID: ${req.user?.id}, Tenant ID: ${tenantId}`);
}
    \end{lstlisting}
    
    \item \textbf{Error Handling Standardizzato}:
    \begin{lstlisting}
// ✅ IMPLEMENTATO
const AppError = require('../utils/AppError');

if (!req.user || !req.user.id || req.user.id === '') {
    if (required) {
        return next(AppError.unauthorized(
            'Tenant context required. Please authenticate.',
            {
                code: 'TENANT_CONTEXT_MISSING',
                category: 'AUTHORIZATION',
                suggestion: 'Assicurati di essere autenticato e che il middleware requireAuth sia applicato prima di tenantContext.',
            }
        ));
    }
    return next();
}
    \end{lstlisting}
\end{enumerate}

\section{UI/UX Impact (Indiretto per Backend)}

\subsection{Velocità di Risposta}

\begin{itemize}
    \item \textbf{Overhead Minimo}: La creazione di \texttt{req.tenantScope} aggiunge < 1ms di latenza
    \item \textbf{Conversione ObjectId}: Operazione O(1), impatto trascurabile
    \item \textbf{Impatto Totale}: Nessun impatto percepibile dall'utente
\end{itemize}

\subsection{Sicurezza dei Dati}

\begin{importantbox}{Impatto Sicurezza}
TenantContext è \textbf{CRITICO} per la sicurezza del sistema:
\begin{itemize}
    \item \textbf{Isolamento Garantito}: Senza \texttt{req.tenantScope}, BaseService non sa chi è l'utente
    \item \textbf{Fail-Fast}: Se utente non autenticato, errore 401 immediato (previene data leak)
    \item \textbf{Validazione Plugin}: Verifica che i modelli abbiano multi-tenancy (previene bug di sviluppo)
    \item \textbf{Prevenzione IDOR}: \texttt{tenantId} viene iniettato server-side, non dal client
\end{itemize}
\end{importantbox}

\subsection{Messaggi di Errore}

\begin{itemize}
    \item \textbf{Standardizzati}: Codice errore \texttt{TENANT\_CONTEXT\_MISSING} per identificazione facile
    \item \textbf{User-Friendly}: Messaggio chiaro: \texttt{"Tenant context required. Please authenticate."}
    \item \textbf{Debug-Friendly}: Se plugin mancante, errore esplicito con nome Model
\end{itemize}

\subsection{Esperienza Sviluppatore}

\begin{itemize}
    \item \textbf{API Consistente}: \texttt{req.tenantScope} disponibile in tutte le routes protette
    \item \textbf{Helper Methods}: \texttt{model()}, \texttt{create()}, \texttt{owns()} rendono il codice più leggibile
    \item \textbf{Documentazione Implicita}: Se conosci TenantContext, capisci come funziona il multi-tenancy
    \item \textbf{Onboarding Veloce}: Nuovi sviluppatori capiscono subito come accedere al contesto tenant
\end{itemize}

\section{Esempi di Utilizzo}

\subsection{Esempio 1: Route Protetta Standard}

\begin{lstlisting}
// routes/api.js
const { requireAuth } = require('../middleware/auth');
const { tenantContext } = require('../middleware/tenantContext');

// Middleware applicato a tutte le routes
router.use(requireAuth);
router.use(tenantContext({ required: true }));

// Route handler
router.get('/projects', async (req, res) => {
    // req.tenantScope è disponibile
    const projects = await projectService.find(req.tenantScope);
    res.json({ success: true, data: projects });
});
\end{lstlisting}

\subsection{Esempio 2: Route Pubblica con Context Opzionale}

\begin{lstlisting}
// Route che funziona sia per utenti autenticati che non
router.use(optionalAuth); // Aggiunge req.user se token presente
router.use(tenantContext({ required: false })); // Non richiede autenticazione

router.get('/public-data', async (req, res) => {
    if (req.tenantScope) {
        // Utente autenticato: mostra dati personalizzati
        const data = await service.findPersonalized(req.tenantScope);
        return res.json({ success: true, data });
    }
    // Utente non autenticato: mostra dati pubblici
    const data = await service.findPublic();
    res.json({ success: true, data });
});
\end{lstlisting}

\subsection{Esempio 3: Utilizzo Helper Methods}

\begin{lstlisting}
// Controller
router.post('/projects', async (req, res) => {
    // Metodo 1: Usando helper create()
    const project = await req.tenantScope.create(Project, {
        name: req.body.name,
        code: req.body.code,
    });
    
    // Metodo 2: Usando helper model()
    const projects = await req.tenantScope.model(Project)
        .find({ status: 'active' })
        .sort('-createdAt')
        .limit(10);
    
    // Metodo 3: Verifica ownership esplicita
    if (!await req.tenantScope.owns(Project, req.params.id)) {
        throw AppError.forbidden('Non puoi accedere a questa risorsa');
    }
    
    // Metodo 4: Filtro base per query complesse
    const stats = await Project.aggregate([
        { $match: { ...req.tenantScope.filter, status: 'active' } },
        { $group: { _id: '$category', total: { $sum: 1 } } }
    ]);
    
    res.json({ success: true, data: project });
});
\end{lstlisting}

\subsection{Esempio 4: withTenantAuth Helper}

\begin{lstlisting}
// routes/goals.js
const { requireAuth } = require('../middleware/auth');
const { withTenantAuth } = require('../middleware/tenantContext');

// Applica auth + tenant context in un'unica chiamata
router.use(...withTenantAuth(requireAuth));

// Equivalente a:
// router.use(requireAuth);
// router.use(tenantContext({ required: true }));
\end{lstlisting}

\section{Pattern e Principi}

\subsection{Middleware Pattern}

\textbf{Definizione}: Middleware è una funzione che ha accesso all'oggetto request, response e alla funzione \texttt{next()} nel ciclo request-response di Express.

\textbf{Implementazione in TenantContext}:
\begin{itemize}
    \item \textbf{Input}: \texttt{req} (con \texttt{req.user} da auth middleware)
    \item \textbf{Output}: \texttt{req.tenantScope} iniettato nella request
    \item \textbf{Chain}: Chiama \texttt{next()} per passare al middleware successivo
\end{itemize}

\textbf{Vantaggi}:
\begin{itemize}
    \item Separazione delle responsabilità (autenticazione vs contesto tenant)
    \item Riusabilità (applicabile a qualsiasi route)
    \item Testabilità (facile mockare \texttt{req.user})
\end{itemize}

\subsection{Dependency Injection (Implicita)}

\textbf{Definizione}: Iniettare dipendenze invece di crearle internamente.

\textbf{Implementazione in TenantContext}:
\begin{itemize}
    \item \texttt{req.tenantScope} viene "iniettato" nella request
    \item I service ricevono \texttt{req.tenantScope} come parametro (non lo creano)
    \item Facilita testing (mock di \texttt{req.tenantScope})
\end{itemize}

\subsection{Principi SOLID}

\begin{enumerate}
    \item \textbf{[S] Single Responsibility}: TenantContext ha UNA sola responsabilità: preparare il contesto tenant
    \item \textbf{[O] Open/Closed}: Aperto per estensione (helper methods), chiuso per modifica
    \item \textbf{[L] Liskov Substitution}: Qualsiasi middleware che crea \texttt{req.tenantScope} può sostituire TenantContext
    \item \textbf{[I] Interface Segregation}: Helper methods granulari (\texttt{model()}, \texttt{create()}, \texttt{owns()})
    \item \textbf{[D] Dependency Inversion}: Dipende da astrazione (\texttt{req.user} interface), non da implementazione
\end{enumerate}

\section{Flusso Completo: Da Request a Database}

\begin{lstlisting}
1. [HTTP Request] 
   GET /api/projects
   Cookie: accessToken=eyJhbGciOiJIUzI1NiIsInR5cCI6IkpXVCJ9...

2. [requireAuth Middleware]
   - Estrae token dal cookie
   - Verifica firma JWT
   - Verifica blacklist
   - Popola req.user = { id: "507f1f77bcf86cd799439011", email: "user@example.com" }

3. [tenantContext Middleware]
   - Legge req.user.id
   - Converte in ObjectId: tenantId = ObjectId("507f1f77bcf86cd799439011")
   - Crea req.tenantScope = { tenantId, model(), create(), owns(), filter }

4. [Route Handler]
   router.get('/projects', async (req, res) => {
       const projects = await projectService.find(req.tenantScope);
   });

5. [BaseService.find()]
   - Estrae tenantId da req.tenantScope.tenantId
   - Crea filtro: { user: ObjectId("507f1f77bcf86cd799439011") }
   - Esegue: Project.find({ user: ObjectId("507f1f77bcf86cd799439011") })

6. [Mongoose Query]
   - Esegue query MongoDB con filtro { user: ObjectId("507f1f77bcf86cd799439011") }
   - Ritorna solo progetti dell'utente autenticato

7. [Response]
   res.json({ success: true, data: projects })
\end{lstlisting}

\section{Conclusioni}

TenantContext è un componente \textbf{critico} che funge da "Guardiano del Cancello" tra autenticazione e accesso ai dati. Senza di lui, BaseService non sa chi è l'utente e il sistema crolla.

\textbf{Punti Chiave}:
\begin{itemize}
    \item ✅ Estrae \texttt{req.user.id} dal middleware di autenticazione
    \item ✅ Converte string ID in ObjectId per compatibilità MongoDB
    \item ✅ Inietta \texttt{req.tenantScope} con helper methods eleganti
    \item ✅ Gestisce edge cases (utente non autenticato, ID malformato)
    \item ✅ Fail-fast con messaggi di errore chiari
    \item ✅ Validazione plugin automatica (previene bug di sviluppo)
    \item ✅ \textbf{Correzioni Applicate}: Validazione ObjectId rigorosa, stringa vuota, Object.freeze(), AppError standardizzato
\end{itemize}

\textbf{Correzioni Implementate}:
\begin{enumerate}
    \item ✅ \textbf{Error Handling Standardizzato}: Uso di \texttt{AppError.unauthorized()} e \texttt{AppError.validation()} invece di \texttt{res.status().json()} manuale
    \item ✅ \textbf{Validazione ObjectId Rigorosa}: Controllo con \texttt{mongoose.Types.ObjectId.isValid()} prima della conversione
    \item ✅ \textbf{Validazione Stringa Vuota}: Check esplicito per \texttt{req.user.id === ''}
    \item ✅ \textbf{Object.freeze()}: Prevenzione modifiche accidentali su \texttt{req.tenantScope}
\end{enumerate}

\textbf{Raccomandazioni Future}:
\begin{enumerate}
    \item Convertire in TypeScript per type safety
    \item Considerare logging opzionale per debug (se \texttt{process.env.DEBUG\_TENANT === 'true'})
\end{enumerate}

\textbf{Relazione con BaseService}:
\begin{itemize}
    \item TenantContext prepara il "carburante" (\texttt{req.tenantScope})
    \item BaseService usa il "carburante" per filtrare query
    \item Senza TenantContext, BaseService non sa chi è l'utente → sistema crolla
    \item Insieme, garantiscono isolamento completo multi-tenant
\end{itemize}

\vspace{1cm}

\textit{Documento generato il \today}

\end{document}
